\section*{الفصل \ref{ch:g10:light-spectra} --- أطياف الضوء ورسالة الضوء}

\begin{solution}{exo:g10:light-spectra:1}
$\qty{254}{nm} < \qty{400}{nm}$: \omterm{def:g10:light-spectra:wavelength}{فوق بنفسجي}. $\qty{480}{nm}$: مرئي،
أزرق. $\qty{589}{nm}$: مرئي، أصفر برتقالي. $\qty{656}{nm}$: مرئي،
أحمر. $\qty{1064}{nm} > \qty{800}{nm}$: \omterm{def:g10:light-spectra:wavelength}{تحت أحمر} --- وهذا ما يجعل ليزر
القطع خطرًا مضاعفًا: قويًّا وغير مرئي.
\end{solution}

\begin{solution}{exo:g10:light-spectra:2}
\emph{1.} شريط مستمر من الألوان، من الأحمر إلى البنفسجي: أي \omterm{def:g10:light-spectra:white}{طيف}
ضوء الشمس.

\emph{2.} الطرف البنفسجي --- فالموشور يحرف أطوال الموجة القصيرة
أكثر من غيرها.

\emph{3.} بقعة خضراء واحدة: منحرفة، لكنها غير منتشرة في شريط.
فضوء الليزر \omterm{def:g10:light-spectra:white}{أحادي اللون} --- طول موجة واحد لا غير، فليس ثمة ما
يُرتَّب.
\end{solution}

\begin{solution}{exo:g10:light-spectra:3}
(أ) مستمر (شعيرة كثيفة ساخنة). (ب) إصدار خطي (غاز مثار
تحت ضغط منخفض). (ج) مستمر (معدن منصهر كثيف). (د) إصدار خطي
(الخط الواحد عند $\qty{589}{nm}$ أساسًا). (ه) امتصاص: أي
\omterm{def:g10:light-spectra:continuous}{الطيف المستمر} للسطح الساخن، تخدشه الخطوط القاتمة
للغلاف الجوي الأبرد للشمس.
\end{solution}

\begin{solution}{exo:g10:light-spectra:4}
\emph{1.} قطرات المطر: فكل قطرة تتلقى ضوء الشمس الأبيض وتحلّله
وتردّه مرتَّبًا بحسب اللون.

\emph{2.} الشمس توفّر \omterm{def:g10:light-spectra:white}{الضوء الأبيض}؛ والمطر يوفّر ملايين
المواشير الصغيرة؛ ولا تستقيم الهندسة إلا للضوء الذي تردّه القطرات
نحو الراصد، ولذلك وجب أن تكون الشمس خلف ظهر الراصد
والمطر أمامه.

\emph{3.} لا تضيف القطرات ضوءًا من عندها --- إنما ترتّب ما
تتلقاه. فإذا عرض ضوء الشمس المرتَّب كل الألوان، فكل لون
كان موجودًا في ضوء الشمس أصلًا: أي أن ضوء الشمس \omterm{def:g10:light-spectra:white}{ضوء أبيض}.
\end{solution}

\begin{solution}{exo:g10:light-spectra:5}
\emph{1.} أحمر باهت $<$ برتقالي زاهٍ $<$ أبيض تقريبًا.

\emph{2.} كلاهما: فالتوهّج يزداد \emph{سطوعًا} (الأثر 1)، ولونه
ينزاح من الأحمر نحو الجهة الزرقاء (الأثر 2، أي تقلّص الذروة
$\lambda_{\max}$)، حتى يضيء من الشريط المرئي جزء أكبر فأكبر إضاءة قوية ---
والشريط الكامل يُقرأ أبيض.

\emph{3.} الجسم المتوهّج بالبياض يصدر إصدارًا قويًّا في الشريط المرئي
كله؛ أما المتوهّج بالحمرة فلا يبلغ إلا طرف أطوال الموجة الطويلة. وبحسب
قانون فين يكون الأول هو الأسخن.
\end{solution}

\begin{solution}{exo:g10:light-spectra:6}
\emph{1.}
\[
T = \frac{\qty{2.90e-3}{m.K}}{\qty{5.00e-7}{m}} = \qty{5800}{K}.
\]

\emph{2.} $\theta = 5800 - 273 = \qty{5527}{\degreeCelsius} \approx
\qty{5500}{\degreeCelsius}$.

\emph{3.} $\qty{400}{nm} < \qty{500}{nm} < \qty{800}{nm}$: فالذروة
مرئية، في الأخضر.
\end{solution}

\begin{solution}{exo:g10:light-spectra:7}
\emph{1.} الشمعة:
$\lambda_{\max} = \qty{2.90e-3}{m.K} / \qty{1800}{K} = \qty{1.61e-6}{m}
= \qty{1.6}{\micro m}$، في \omterm{def:g10:light-spectra:wavelength}{تحت الأحمر}. صفيحة التسخين:
$\lambda_{\max} = \qty{2.90e-3}{m.K} / \qty{500}{K} = \qty{5.8e-6}{m}
= \qty{5.8}{\micro m}$، \omterm{def:g10:light-spectra:wavelength}{تحت أحمر} بعيد.

\emph{2.} عند $\qty{500}{K}$ يكون كل الإشعاع تقريبًا \omterm{def:g10:light-spectra:wavelength}{تحت أحمر}:
فالجلد يمتصّه ويحسّ الحرارة، بينما لا تتلقى العين شيئًا --- ساخنة
وسوداء في آن.

\emph{3.} \omterm{def:g10:light-spectra:continuous}{الطيف المستمر} عريض: فحتى وذروته عند
$\qty{1.6}{\micro m}$، يبقى \omterm{def:g10:light-spectra:white}{لطيف} الشمعة ذيل مرئي عند
الطرف الأحمر الأصفر --- ضعيف، لكنه يكفي العين المتكيّفة مع الظلام.
\end{solution}

\begin{solution}{exo:g10:light-spectra:8}
المصباح A: الخطوط $434$ و$486$ و$\qty{656}{nm}$ توافق الهيدروجين، والخط
البنفسجي الخافت يكمّل القائمة عند $\qty{410}{nm}$: فهو الهيدروجين.

المصباح B: يقع الخط عند $\qty{589}{nm}$، وهو خط الصوديوم المرئي القوي
الوحيد. ولا يُستبعد الهيليوم بالخط نفسه (فللهيليوم فعلًا خط
عند $\qty{588}{nm}$ القريب) بل بالخطوط \emph{الغائبة}: إذ كان الهيليوم
سيُظهر كذلك $447$ و$502$ و$\qty{668}{nm}$، ولا يظهر منها شيء. فالعنصر
يُعيَّن بقائمته الكاملة: أي أن المصباح B يحوي صوديومًا.
\end{solution}

\begin{solution}{exo:g10:light-spectra:9}
\emph{1.} قوس قزح مستمر يخدشه خط قاتم رفيع واحد: أي
\omterm{def:g10:light-spectra:absorption}{طيف امتصاص} الصوديوم.

\emph{2.} عند $\qty{589}{nm}$: فالغاز يمتصّ أطوال الموجة التي يقدر
على إصدارها بالضبط (\cref{prop:g10:light-spectra:kirchhoff})، وخط الصوديوم
المرئي القوي عند $\qty{589}{nm}$.

\emph{3.} \omterm{def:g10:light-spectra:white}{طيف} إصدار الصوديوم: خط واحد ساطع أصفر برتقالي
عند $\qty{589}{nm}$ على خلفية سوداء --- \omterm{def:g10:light-spectra:wavelength}{طول الموجة} نفسه،
مع تبادل الساطع والقاتم بين الطيفين.
\end{solution}

\begin{solution}{exo:g10:light-spectra:10}
\emph{1.} $410$ و$434$ و$486$ و$\qty{656}{nm}$: أي قائمة الهيدروجين
كاملة. و$517$ و$\qty{518}{nm}$: زوج المغنيزيوم. فالعنصران كلاهما
حاضر.

\emph{2.} في الغلاف الجوي للشمس، أي الغاز الأبرد الأرقّ فوق
السطح المتوهّج. فالسطح يرسل طيفًا مستمرًّا؛ وينزع
الغلاف الجوي أطوال موجته الخاصة من الضوء العابر، فتظهر
الخطوط ضوءًا مفقودًا --- قاتمة على الخلفية الساطعة.
\end{solution}

\begin{solution}{exo:g10:light-spectra:11}
\emph{1.} $T = 37 + 273 = \qty{310}{K}$، ومنه
\[
\lambda_{\max} = \frac{\qty{2.90e-3}{m.K}}{\qty{310}{K}}
= \qty{9.4e-6}{m} = \qty{9.4}{\micro m}.
\]

\emph{2.} \omterm{def:g10:light-spectra:wavelength}{تحت أحمر} بعيد --- أي أكثر من عشرة أضعاف طول موجة الضوء
الأحمر.

\emph{3.} عند $\qty{310}{K}$ يكون الجزء المرئي من الإشعاع
مهمَلًا تمامًا: فلا توهّج مرئي مهما أظلمت الغرفة. أما الكاميرا
الحرارية فتحمل حسّاسًا مضبوطًا على الإشعاع المصدَر فعلًا، قرب
$\qty{9}{\micro m}$، فترى الأجسام الدافئة تسطع في مواجهة محيط
أبرد.
\end{solution}

\begin{solution}{exo:g10:light-spectra:12}
\emph{1.} $\lambda_{\max} = \qty{2.90e-3}{m.K} / \qty{2700}{K}
= \qty{1.07e-6}{m} \approx \qty{1.1}{\micro m}$: \omterm{def:g10:light-spectra:wavelength}{تحت أحمر} قريب.

\emph{2.} تقع ذروة \omterm{def:g10:light-spectra:white}{الطيف} وراء المرئي: فمعظم الاستطاعة
الكهربائية يعود \omterm{def:g10:light-spectra:wavelength}{تحت أحمر} غير مرئي --- أي حرارة --- ولا يضيء الغرفة
إلا ذيل مرئي صغير. فالمصباح، مصدرَ إضاءة، يبدّد معظم
استطاعته؛ أما مدفأةً فهو شبه مثالي.

\emph{3.} $T = \qty{2.90e-3}{m.K} / \qty{5.50e-7}{m} \approx
\qty{5300}{K}$ --- أي أعلى بكثير من درجة انصهار التنغستن البالغة
$\qty{3700}{K}$. فلا شعيرة صلبة يمكن دفعها إلى درجات حرارة
ضوء النهار، فاضطرت مصابيح شبيهة بضوء النهار إلى هجر التوهّج
بالكلية (الأنابيب الفلورية، وثنائيات LED)، فأنتجت ضوءًا مرئيًّا دون
تسخين جسم إلى $\qty{5300}{K}$.
\end{solution}

\begin{solution}{exo:g10:light-spectra:13}
\emph{1.} $T = \qty{2.90e-3}{m.K} / \qty{4.20e-7}{m} \approx
\qty{6900}{K}$: أسخن من الشمس عند $\qty{5800}{K}$.

\emph{2.} الهيدروجين (فالمجموعة كاملة $410$ و$434$ و$486$ و$\qty{656}{nm}$
حاضرة) والكالسيوم (خطاه كلاهما، $393$ و$\qty{397}{nm}$).

\emph{3.} لا. فخط الصوديوم عند $\qty{589}{nm}$ غائب، ولو كان الصوديوم
وفيرًا في الغلاف الجوي لطبعه حتمًا على الخلفية
المستمرة. فإن كان ثمة صوديوم أصلًا، فآثار خافتة لا
تُكتشف.
\end{solution}

\begin{solution}{exo:g10:light-spectra:14}
\emph{1.} الهيدروجين يفسّر $410$ و$434$ و$486$ و$\qty{656}{nm}$؛
والهيليوم يفسّر $447$ و$502$ و$588$ و$\qty{668}{nm}$. فكل خط مرصود
مفسَّر: أي أن المزيج هيدروجين وهيليوم.

\emph{2.} الهيليوم تفسير \emph{كافٍ} لأن خطوطه الثلاثة
الأخرى ($447$ و$502$ و$\qty{668}{nm}$) كلها حاضرة --- فيكون الخط
$\qty{588}{nm}$ متوقَّعًا من الهيليوم وحده، وليس في
\omterm{def:g10:light-spectra:white}{الطيف} ما \emph{يقتضي} صوديومًا. أما حسم المسألة فيحتاج
تحليلًا لا جدلًا: فمِطياف أدقّ يفصل خط الهيليوم
الواحد عند $\qty{587.6}{nm}$ عن خط الصوديوم، وهو في الحقيقة زوج متقارب عند
$589.0$ و$\qty{589.6}{nm}$. فرؤية خط واحد هناك تبرّئ الصوديوم؛
ورؤية ثلاثة تدينه.
\end{solution}

\begin{solution}{exo:g10:light-spectra:15}
\emph{1.} إبط الجوزاء:
$T = \qty{2.90e-3}{m.K} / \qty{8.5e-7}{m} \approx \qty{3400}{K}$. رجل الجبار:
$T = \qty{2.90e-3}{m.K} / \qty{2.4e-7}{m} \approx \qty{12000}{K}$.

\emph{2.} بما أن $\lambda_{\max} \times T$ هو الثابت نفسه في الحالتين، فإن
$\dfrac{T_{\text{رجل الجبار}}}{T_{\text{إبط الجوزاء}}}
= \dfrac{\qty{850}{nm}}{\qty{240}{nm}} \approx 3.5$.

\emph{3.} تقع ذروة إبط الجوزاء في \omterm{def:g10:light-spectra:wavelength}{تحت الأحمر}: فطيفه داخل الشريط
المرئي أقوى ما يكون عند الطرف الأحمر، فيبدو النجم أحمر. أما رجل الجبار
فذروته في \omterm{def:g10:light-spectra:wavelength}{فوق البنفسجي}: فشريطه المرئي أقوى ما يكون عند الطرف الأزرق، ومنه
لونه الأبيض المزرقّ.

\emph{4.} يغطي \omterm{def:g10:light-spectra:continuous}{الطيف المستمر} مجالًا هائلًا من أطوال الموجة. فحتى
وذروته خارج الشريط المرئي، يصبّ النجم ضوءًا وفيرًا
\emph{داخل} ذلك الشريط --- فالذروة تقرّر الصبغة لا الرؤية.
\end{solution}

\begin{solution}{pb:g10:light-spectra:1}
\textbf{1.} من نحو $\qty{400}{nm}$ (الطرف البنفسجي) إلى نحو
$\qty{800}{nm}$ (الطرف الأحمر)؛ ويقع $\qty{550}{nm}$ قرب الوسط، في
الأخضر.

\textbf{2.} $\qty{410}{nm}$: بنفسجي؛ $\qty{434}{nm}$: بنفسجي مزرقّ؛
$\qty{486}{nm}$: أزرق مخضرّ؛ $\qty{656}{nm}$: أحمر.

\textbf{3.} توهّج \omterm{def:g10:light-spectra:white}{أحادي اللون} أصفر برتقالي. ولون السيارة هو
مزيج أطوال الموجة التي يعكسها طلاؤها؛ وتحت مصباح لا يعرض
\emph{إلا} $\qty{589}{nm}$، لا يقدر أي طلاء إلا على عكس قدر أكبر أو أصغر
من \omterm{def:g10:light-spectra:wavelength}{طول الموجة} الواحد ذاك --- فتبدو السيارات كلها بدرجات من البرتقالي
والرمادي.

\textbf{4.} يصدر كل عنصر قائمة ثابتة من الخطوط، ولا يشترك عنصران
في القائمة نفسها
(\cref{prop:g10:light-spectra:fingerprint}): فمجموعة خطوط كاملة
تعيّن إذن العنصر كما تعيّنه البصمة. أما التوافق
\emph{الجزئي} فلا يبرهن على كثير --- إذ يمكن أن يملك عنصران متمايزان
خطين مفردين شبه متطابقين (الهيليوم عند $\qty{588}{nm}$، والصوديوم عند
$\qty{589}{nm}$) --- ولذلك يقتضي ادّعاء الحضور كل الخطوط القوية
في القائمة.

\textbf{5.} (أ) قوس قزح مستمر فيه خط قاتم عند
$\qty{589}{nm}$: أي \omterm{def:g10:light-spectra:absorption}{طيف الامتصاص}. (ب) خط ساطع واحد عند
$\qty{589}{nm}$ على سواد: أي \omterm{def:g10:light-spectra:emission}{طيف الإصدار} --- \omterm{def:g10:light-spectra:wavelength}{طول الموجة} نفسه،
مع تبادل الساطع والقاتم
(\cref{prop:g10:light-spectra:kirchhoff}).

\textbf{6.} يزداد \omterm{def:g10:light-spectra:white}{الطيف} سطوعًا عند كل طول موجة، وتنتقل ذروته
$\lambda_{\max}$ نحو أطوال الموجة الأقصر، وفق
$\lambda_{\max} \times T = \qty{2.90e-3}{m.K}$
(\cref{prop:g10:light-spectra:wien}).

\textbf{7.} $\lambda_{\max} = \qty{2.90e-3}{m.K} / \qty{3400}{K}
= \qty{8.5e-7}{m} = \qty{850}{nm}$: \omterm{def:g10:light-spectra:wavelength}{تحت أحمر} قريب، بعد الحافة الحمراء
بقليل --- ولهذا يكون ضوء الهالوجين أدفأ صبغةً من ضوء النهار.

\textbf{8.} $T = \qty{2.90e-3}{m.K} / \qty{5.00e-7}{m} = \qty{5800}{K}$.

\textbf{9.} الذروة لطيفة لا سنّية: فالشمس تصدر الشريط المرئي
كله بقوة متقاربة. ومزج شريط ألوان كامل معًا
يُقرأ أبيض --- فتبدو الشمس بيضاء (مصفرّة عبر غلافنا
الجوي)، ولا تبدو خضراء أبدًا.

\textbf{10.} القضيب البارد: ذروته عميقًا في \omterm{def:g10:light-spectra:wavelength}{تحت الأحمر}، ولا ذيل مرئي ---
فيُحسّ الإشعاع ولا يُرى شيء. وإذا سخن: دفع \omterm{def:g10:light-spectra:white}{الطيف} المتزايد سطوعًا
أول ذيل مرئي وراء $\qty{800}{nm}$ --- فيضيء الطرف الأحمر
وحده: أحمر باهت. وإذا ازداد سخونة: امتلأ الشريط المرئي كله بينما
يزداد كل شيء سطوعًا: برتقالي، ثم نحو الأبيض.

\textbf{11.} تأتي الخلفية المستمرة من السطح المتوهّج الكثيف
الساخن؛ أما الخطوط القاتمة فيطبعها الغلاف الجوي الأبرد الأرقّ
فوقه، وهو يمتصّ أطوال موجته الخاصة من الضوء العابر
خلاله.

\textbf{12.} الهيدروجين ($410$ و$434$ و$486$ و$\qty{656}{nm}$: كاملة)،
والصوديوم ($\qty{589}{nm}$)، والكالسيوم ($393$ و$\qty{397}{nm}$: كلاهما)
--- فالعناصر الثلاثة كلها حاضرة في الغلاف الجوي للشمس.

\textbf{13.} أن الهيدروجين يهيمن على الغلاف الجوي. وتوافق القياسات
الحديثة على ذلك وتحدّده عدديًّا: فالشمس ثلاثة أرباعها تقريبًا
هيدروجين بالكتلة، وأكثر ما بقي هيليوم.

\textbf{14.} لم يوافق الخط أي عنصر معروف؛ وبما أن قائمة كل عنصر
ثابتة وفريدة، فالخط الذي لا ينتمي إلى أي قائمة لا بد أن ينتمي
إلى عنصر لم يُرَ قط في مختبر. فأُعلن عن عنصر جديد
وسُمّي الهيليوم، نسبةً إلى \emph{هيليوس}، أي الشمس. وقد
عُزل على الأرض سنة 1895 --- فتأكد التنبؤ بعد سبع وعشرين
سنة.

\textbf{15.} في مواجهة الخلفية المستمرة الباهرة، \emph{ينزع}
الغلاف الجوي الضوء عند أطوال موجته الخاصة: فتكون الخطوط قاتمة.
وخلال الكسوف يحجب القمر الخلفية؛ فيقف الغلاف الجوي الرقيق
عندئذٍ وحده أمام سماء مظلمة، ولا يُرى إلا ضوؤه المعاد إصداره:
أطوال الموجة نفسها، ساطعة الآن. غاز واحد وقائمة واحدة
--- والإشارة تتوقف على ما يقف خلفه.

\textbf{16.} $T = \qty{2.90e-3}{m.K} / \qty{2.90e-7}{m}
= \qty{10000}{K}$.

\textbf{17.} أبيض مزرقّ: \omterm{def:g10:light-spectra:white}{فالطيف} داخل الشريط المرئي أقوى ما يكون
عند الطرف الأزرق. ولا خفاء البتة: \omterm{def:g10:light-spectra:continuous}{فالطيف المستمر}
يغمر الشريط المرئي كله في طريقه إلى الذروة فوق البنفسجية ---
فالذروة تثبّت الصبغة لا السطوع.

\textbf{18.} الخطوط القاتمة $410$ و$434$ و$486$ و$\qty{656}{nm}$ هي
بصمة الهيدروجين كاملة: فالهيدروجين يهيمن على غلاف النسر الواقع.

\textbf{19.} $T_{\text{النسر الواقع}} / T_{\text{الشمس}} = 10000 / 5800 \approx
1.7$: فالنسر الواقع أسخن، ومن ثم أكثر زرقة (ذروته عند $\qty{290}{nm}$ مقابل
$\qty{500}{nm}$)، وهو، بحسب الجزء الأول من
\cref{prop:g10:light-spectra:wien}، أشدّ سطوعًا لكل \omterm{def:g10:orders-of-magnitude:unit}{وحدة} من السطح المتوهّج
عند كل طول موجة.

\textbf{20.} بطاقة هوية النسر الواقع --- درجة حرارة السطح: نحو
$\qty{10000}{K}$؛ العنصر المهيمن في الغلاف الجوي: الهيدروجين؛ اللون
الظاهري: أبيض مزرقّ. وكل ذلك مقروء من حزمة ضوء واحدة، بموشور
وقائمة خطوط مخبرية وعملية قسمة واحدة. أما ``أبدًا'' كونت فقد
صدرت سنة 1835 وانقضت سنة 1859: أي أنها دامت أربعًا وعشرين سنة.
\end{solution}
